\documentclass[11pt,a4paper]{article}
\usepackage[margin=2.5cm]{geometry}
\usepackage[T1]{fontenc}
\usepackage{amsmath,amssymb,amsthm}
\usepackage{enumitem}
\usepackage{fancyhdr}
\usepackage[hidelinks]{hyperref}

% Replace these four lines with your details; they fill the header and the title.
\newcommand{\course}{MATH 221: Linear Algebra}
\newcommand{\assignment}{Problem Set 4}
\newcommand{\student}{Jonas Weber (ID 20418733)}
\newcommand{\duedate}{Due 20 March 2026}

\pagestyle{fancy}
\fancyhf{}
\lhead{\course}
\chead{\assignment}
\rhead{\student}
\cfoot{\thepage}
\setlength{\headheight}{14pt}
\setlength{\parindent}{0pt}
\setlength{\parskip}{0.6em}

% \problem{points} starts a numbered problem; the solution environment holds your answer.
\newcounter{problem}
\newcommand{\problem}[1]{\stepcounter{problem}\bigskip\noindent\textbf{Problem \theproblem} \textnormal{(#1 points)}\par\nopagebreak}
\newenvironment{solution}{\par\noindent\textit{Solution.}\ }{\hfill$\square$\par}

\begin{document}

\begin{center}
  {\Large\bfseries \course}\\[4pt]
  {\large \assignment}\\[4pt]
  \student \quad\textbar\quad \duedate\\[2pt]
  Instructor: Dr. Elena Marchetti
\end{center}
\hrule

\problem{10}
Solve the system
\[
\begin{aligned}
  x + 2y - z &= 3,\\
  2x - y + 3z &= 7,\\
  -x + 3y + 2z &= 0.
\end{aligned}
\]

\begin{solution}
Write the augmented matrix and reduce it. Subtracting twice the first row from the second, and adding the first row to the third, gives
\[
\left[\begin{array}{rrr|r}
  1 & 2 & -1 & 3\\
  2 & -1 & 3 & 7\\
  -1 & 3 & 2 & 0
\end{array}\right]
\longrightarrow
\left[\begin{array}{rrr|r}
  1 & 2 & -1 & 3\\
  0 & -5 & 5 & 1\\
  0 & 5 & 1 & 3
\end{array}\right].
\]
Adding the second row to the third gives $6z = 4$, so $z = \tfrac{2}{3}$. The second row then gives $-5y + 5z = 1$, so $y = z - \tfrac{1}{5} = \tfrac{7}{15}$, and the first row gives $x = 3 - 2y + z = \tfrac{41}{15}$. Substituting into the second equation confirms the answer: $\tfrac{82}{15} - \tfrac{7}{15} + \tfrac{30}{15} = 7$.
\end{solution}

\problem{10}
Let $A = \begin{pmatrix} 4 & 1 \\ 2 & 3 \end{pmatrix}$. Find the eigenvalues of $A$ and an eigenvector for each.

\begin{solution}
The characteristic polynomial is
\begin{equation}\label{eq:char}
  \det(A - \lambda I) = (4 - \lambda)(3 - \lambda) - 2 = \lambda^2 - 7\lambda + 10 = (\lambda - 2)(\lambda - 5),
\end{equation}
so the eigenvalues are $\lambda_1 = 2$ and $\lambda_2 = 5$.

For $\lambda_1 = 2$ we solve $(A - 2I)\mathbf{v} = \mathbf{0}$, that is $2v_1 + v_2 = 0$, which gives $\mathbf{v}_1 = (1, -2)^{\mathsf{T}}$. For $\lambda_2 = 5$ we solve $-v_1 + v_2 = 0$, which gives $\mathbf{v}_2 = (1, 1)^{\mathsf{T}}$. As a check, the eigenvalues in \eqref{eq:char} sum to $7 = \operatorname{tr} A$ and multiply to $10 = \det A$.
\end{solution}

\problem{15}
Let $V$ be a vector space and let $U$ and $W$ be subspaces of $V$.
\begin{enumerate}[label=(\alph*)]
  \item Prove that $U \cap W$ is a subspace of $V$.
  \item Give an example showing that $U \cup W$ need not be a subspace.
\end{enumerate}

\begin{solution}
\begin{enumerate}[label=(\alph*)]
  \item The zero vector lies in both $U$ and $W$, so $U \cap W$ is not empty. Let $\mathbf{u}, \mathbf{v} \in U \cap W$ and let $c$ be a scalar. Because $U$ is a subspace, $\mathbf{u} + \mathbf{v} \in U$ and $c\mathbf{u} \in U$; the same holds for $W$. Hence $\mathbf{u} + \mathbf{v}$ and $c\mathbf{u}$ lie in $U \cap W$, which is therefore closed under addition and scalar multiplication.
  \item In $V = \mathbb{R}^2$, take $U = \{(x, 0) : x \in \mathbb{R}\}$ and $W = \{(0, y) : y \in \mathbb{R}\}$. Then $(1, 0)$ and $(0, 1)$ are in $U \cup W$, but their sum $(1, 1)$ is in neither $U$ nor $W$.
\end{enumerate}
\end{solution}

\problem{15}
Let $T \colon \mathbb{R}^3 \to \mathbb{R}^2$ be the linear map $T(x, y, z) = (x + y,\; y - z)$. Find a basis of the kernel of $T$ and verify the rank--nullity theorem.

\begin{solution}
A vector $(x, y, z)$ is in the kernel when $x + y = 0$ and $y - z = 0$, that is $x = -y$ and $z = y$. So
\[
  \ker T = \{\, t\,(-1, 1, 1) : t \in \mathbb{R} \,\},
\]
with basis $\{(-1, 1, 1)\}$ and dimension 1. The images $T(1, 0, 0) = (1, 0)$ and $T(0, 0, 1) = (0, -1)$ are linearly independent, so the image is all of $\mathbb{R}^2$ and the rank is 2. The theorem holds: $\operatorname{rank} T + \dim \ker T = 2 + 1 = 3 = \dim \mathbb{R}^3$.
\end{solution}

\bigskip
\hrule
\smallskip
% Many instructors ask for this line; delete it if yours does not.
\textit{Collaboration: I discussed Problem 3 with Hana Suzuki. All written solutions are my own.}

\end{document}
